\documentclass[11pt]{article}
\usepackage[a4paper,margin=18mm]{geometry}
\usepackage[no-math]{fontspec}
\setmainfont{Latin Modern Roman}
\usepackage{amsmath,amssymb,amsthm,xfrac,xspace,hyperref,graphicx,comment}
\excludecomment{DefTralics}
\providecommand{\backmatter}{}
\providecommand{\loccit}{\emph{loc. cit.}}

\providecommand{\bibliofont}{\small}
\newcommand{\ie}{i.e.,\xspace}
\newcommand{\eg}{e.g.,\xspace}
\newcommand{\etc}{etc.\xspace}
\newcommand{\cf}{cf.\xspace}
\newcommand{\resp}{resp.\xspace}
\newcommand{\smaller}{\small}
\newcommand{\Subsection}{\subsection}
\newcommand{\Subsubsection}{\subsubsection}
\newcommand{\skpt}{\smallskip}
\emergencystretch=3em

% Begin literal retained input author-preamble.tex
\usepackage{mathrsfs}
\let\scr\mathscr

\newcommand\mto{\mathchoice{\longmapsto}{\mapsto}{\mapsto}{\mapsto}}
\newcommand\hto{\mathrel{\lhook\joinrel\to}}
\newcommand{\To}[1]{\mathchoice{\xrightarrow{\ts\kern4pt#1\kern3pt}}{\stackrel{#1}{\longrightarrow}}{}{}}
\let\ra\rightarrow

\let\vref\ref
\let\ts\relax
\newcommand{\sbullet}{{\scriptscriptstyle\bullet}}
\newcommand{\ab}{\mathrm{ab}}
\newcommand{\an}{\mathrm{an}}

\newenvironment{enumeraten}
{\begin{enumerate}\setcounter{enumi}{\value{equation}}\itemindent0pt}
{\setcounter{equation}{\value{enumi}}\end{enumerate}}

\usepackage{mathtools}
\usepackage{tikz}
\usepackage{tikz-cd}

\tikzcdset{arrow style=Latin Modern}
\newcommand{\dradpl}{{\hspace*{-2.5mm}\begin{tikzcd}[ampersand replacement=\&,column sep = scriptsize]\mbox{}\ar[r,dashed]\&\mbox{}\end{tikzcd}\hspace*{-2.5mm}}}
\newcommand{\dratst}{{\hspace*{-2.5mm}\begin{tikzcd}[ampersand replacement=\&,column sep = small]\mbox{}\ar[r,dashed]\&\mbox{}\end{tikzcd}\hspace*{-2mm}}}
\let\drascst\dratst
\newcommand{\dra}{\mathchoice{\dradpl}{\dratst}{\drascst}{\drascst}}
\let\dasharrow\dra

\theoremstyle{plain}
\newtheorem{thm}{Theorem}[section]
\newtheorem{claim}[thm]{Claim}
\newtheorem{consequence}[thm]{Consequence}
\newtheorem{fact}[thm]{Fact}
\newtheorem{lem}[thm]{Lemma}
\newtheorem{prop}[thm]{Proposition}

\theoremstyle{definition}
\newtheorem{defn}[thm]{Definition}
\newtheorem{asswlog}[thm]{Assumption w.l.o.g.}
\newtheorem{construction}[thm]{Construction}
\newtheorem{notation}[thm]{Notation}
\newtheorem{rem}[thm]{Remark}

\numberwithin{equation}{thm}
\numberwithin{enumi}{thm}

\DeclareMathOperator{\Aut}{Aut}
\DeclareMathOperator{\codim}{codim}
\DeclareMathOperator{\img}{img}
\DeclareMathOperator{\Pic}{Pic}
\DeclareMathOperator{\rank}{rank}
\newcommand{\reg}{\mathrm{reg}}
\DeclareMathOperator{\sEnd}{\sE\negthinspace \mathit{nd}}
\DeclareMathOperator{\sing}{sing}
\DeclareMathOperator{\Sym}{Sym}
\DeclareMathOperator{\supp}{supp}
\DeclareMathOperator{\tor}{tor}

\DeclareMathOperator{\alb}{alb}
\DeclareMathOperator{\Alb}{Alb}
\DeclareMathOperator{\diff}{d}
\DeclareMathOperator{\Div}{Div}
\DeclareMathOperator{\Gal}{Gal}
\DeclareMathOperator{\GL}{GL}
\DeclareMathOperator{\iit}{iit}
\DeclareMathOperator{\Iit}{Iit}
\DeclareMathOperator{\Lie}{Lie}
\DeclareMathOperator{\PGL}{\bP\negthinspace\GL}
\DeclareMathOperator{\Rad}{Rad}
\DeclareMathOperator{\sh}{sha}
\DeclareMathOperator{\Sh}{Sha}
\DeclareMathOperator{\SL}{SL}

\newcommand{\sA}{\scr{A}}
\newcommand{\sE}{\scr{E}}
\newcommand{\sF}{\scr{F}}
\newcommand{\sL}{\scr{L}}
\newcommand{\sN}{\scr{N}}
\newcommand{\sO}{\scr{O}}
\newcommand{\sS}{\scr{S}}
\newcommand{\sT}{\scr{T}}
\newcommand{\sU}{\scr{U}}

\newcommand{\bC}{\mathbb{C}}
\newcommand{\bN}{\mathbb{N}}
\newcommand{\bP}{\mathbb{P}}
\newcommand{\bQ}{\mathbb{Q}}
\newcommand{\bZ}{\mathbb{Z}}

\let\wtilde\widetilde
\let\what\widehat

\hyphenation{com-po-nents}
\hyphenation{pos-i-tive}
\hyphenation{Theo-rem}
\hyphenation{Vojta}

\newcommand{\factor}[2]{\left. \raise 2pt\hbox{$#1$} \right/\hskip -2pt\raise -2pt\hbox{$#2$}}

% End literal retained input author-preamble.tex

\begin{document}
\begin{center}{\Large Stable item 2057}\end{center}
\paragraph{Source and attribution.}
Daniel Greb, Stefan Kebekus, Thomas Peternell, \emph{Projectively flat klt varieties}, Journal de l’École polytechnique — Mathématiques 8 (2021), publisher version, DOI 10.5802/jep.164. \href{https://jep.centre-mersenne.org/articles/10.5802/jep.164/}{Publisher article}. Authors' text: \href{https://creativecommons.org/licenses/by/4.0/}{CC BY 4.0}.
\paragraph{Editorial problem boundary.}
Prove only Proposition 4.1 below for a normal complex space with klt singularities whose reflexive cotangent sheaf is the direct sum of \(\dim X\) copies of one rank-one reflexive sheaf. The conclusion is that all singularities are isolated cyclic quotient singularities. Exact original conventions and reflexive-sheaf definitions, full Section4 with the complete local product lemma and index-one-cover/tensor-power proof are retained. Global projectively flat classification and quasi-Abelian characterization are not selected claims.
\paragraph{Difficulty (editorial): H2.}
The cited klt Lipman--Zariski theorem can make an index-one cover smooth, but it does not show that the resulting cyclic quotient singularities are isolated. The author first proves the local product description along a positive-dimensional singular component, then uses the equal rank-one splitting to force a reflexive tensor power to be locally free. The incompatible product summand supplies the remaining nonroutine isolation argument.
\paragraph{Editorial transcription note.}
Original author language, mathematics and ordering are preserved. Publisher front matter is replaced by this independent article wrapper. Bibliography is recovered from matching cached publisher HTML, including its TeX formulas. No new proof or mathematical repair is supplied. Surrounding original results are retained for conservative context and are not extra selected corpus claims.
\clearpage
\setcounter{section}{1}
% Begin literal retained input author-context.tex
\section{Conventions, notation, and variations of standard facts}

\subsection{Global conventions}

Throughout this paper, all schemes, varieties and morphisms will be defined over
the complex number field. We follow the notation and conventions of
Hartshorne's book \cite{Ha77}. In particular, varieties are always assumed to
be irreducible. For all notation around Mori theory, such as klt spaces and klt
pairs, we refer the reader to \cite{KM98}.

\subsection{Varieties and complex spaces}

In order to keep notation simple, we will sometimes, when there is no danger of
confusion, not distinguish between algebraic varieties and their underlying
complex spaces. Along these lines, if $X$ is a quasi-projective complex
variety, we write $\pi_1(X)$ for the fundamental group of the associated complex
space.

\subsection{Reflexive sheaves}

As in most other papers on the subject, we will frequently consider reflexive
sheaves and take reflexive hulls. Given a normal, quasi-projective variety (or
normal, irreducible complex space) $X$, we write
$\Omega^{[p]}_X := \bigl(\Omega^p_X \bigr)^{**}$ and refer to this sheaf as \emph{the
sheaf of reflexive differentials}. More generally, given any coherent sheaf
$\sE$ on $X$, write $\sE^{[\otimes m]} := \bigl(\sE^{\otimes m} \bigr)^{**}$ and
$\det \sE := \bigl(\wedge^{\rank \sE} \sE \bigr)^{**}$. Given any morphism $f : Y \to X$ of
normal, quasi-projective varieties (or normal, irreducible, complex spaces), we
write $f^{[*]} \sE := (f^* \sE)^{**}$.


% End literal retained input author-context.tex
\setcounter{section}{3}
% Begin literal retained input author-body.tex
\section{Varieties with splitting cotangent sheaves}
\label{sec:4}

One relevant case that keeps appearing in the discussion of varieties with
projectively flat sheaf of reflexive differentials is the one where $\Omega^{[1]}_X$
decomposes as a direct sum of Weil divisorial sheaves. We discuss settings
where there are local or global decompositions of this form. The results of
this section will be used later in the proof of our main result, but might be of
independent interest. In relation to \cite{MR3030068}, the local results are
new, prompted by the appearance of singularities in our context, and the results
on varieties with globally split cotangent sheaf allow us to replace some of the
more involved arguments of \loccit\ by a clearer structure result.

\subsection{Varieties where $\Omega^{[1]}_X$ is locally split}
\label{sect:split}

Given a variety $X$ where $\Omega^{[1]}_X$ is projectively flat, for instance because
the flatness criterion of \cite[Th.\,1.6]{GKP20a} applies, the local description
of projectively flat sheaves, \cite[Prop.\,3.11]{GKP20a}, can often be used to
reduce to the case where $\Omega^{[1]}_X$ decomposes locally. We will show that this
already forces $X$ to have quotient singularities.

\begin{prop}\label{prop:icq}
Let $X$ be a normal complex space with klt singularities. Assume that the
sheaf $\Omega^{[1]}_X$ of reflexive differentials is of the form
\begin{equation}\label{eq:icq}
\Omega^{[1]}_X \cong \sL^{\oplus \dim X},
\end{equation}
where $\sL$ is reflexive of rank one. Then, $X$ has only isolated, cyclic
quotient singularities.
\end{prop}

We prove Proposition~\ref{prop:icq} in Section~\vref{ssec:popicq}.

\begin{rem}
If $X$ is a quasi-projective variety with klt singularities where $\Omega^{[1]}_X$
is of the form $\sL^{\oplus \dim X}$ for a Weil-divisorial sheaf $\sL$, then the proof
of Proposition~\ref{prop:icq} applies verbatim to show that $X$ is
Zariski-locally a quotient of a smooth variety. We do not use this fact
however.
\end{rem}

\subsection{Varieties where $\Omega^{[1]}_X$ is globally split}
\label{sec:4-2}

There are settings where $\Omega^{[1]}_X$ decomposes globally, for instance because
it is projectively flat and $X$ is known to be simply connected. This obviously
puts strong conditions on the global geometry of $X$; see \cite{MR1758581} for
further results in the smooth setup.

The following proposition, including the proof, is due to the referee. He also
suggested Theorem~\ref{thm:druel}, which is stronger than our original result
that included a pseudo-effectivity assumption. Our previous arguments were of
analytic nature and used a theorem of Demailly~\cite{Dem02}.

\begin{prop}\label{prop:druel}
Let $U$ be a connected complex manifold of dimension $n \geq 2$. Assume that
$\Omega^1_U \cong \sL^{\oplus n}$, where $\sL$ is invertible. Then,
$c_1(\sL) = 0 \in H^1 \bigl(U,\, \Omega^1_U \bigr)$.
\end{prop}
\begin{proof}
Consider the induced decomposition on cohomology
\[
H^1 \bigl(U,\, \Omega^1_U\bigr) \cong H^1\bigl(U,\, \sL \bigr)^{\oplus n}.
\]
For reasons that
will become apparent in a second, we are particularly interested in the
coordinate hyperplanes
\[
Z_i := \left\{ (\alpha_1, \dots, \alpha_n) \in H^1\bigl(U,\, \sL \bigr)^{\oplus n} \bigm| \alpha_i = 0
\right\} \subsetneq H^1 \bigl(U,\, \Omega^1_U \bigr).
\]
To see how the $Z_i$ come into play, restate the assumption as a decomposition
of the tangent bundle: $\sT_U \cong \sA_1 \oplus \cdots \oplus \sA_n$. The summands $\sA_\sbullet \cong \sL^*$ are
regular rank-one foliations on $U$, with normal bundles
\[
\sN_i = \sA_1 \oplus \cdots \oplus \what{\sA_i} \oplus \cdots \oplus \sA_n.
\]
In this setting, a theorem of Baum-Bott, \cite[Prop.\,3.3 \&
Cor.~3.4]{MR0261635} but see also \cite[Lem.\,3.1]{MR1760872}, describes the
first Chern class $c_1(\sN_i)$ as lying in the subspace
\[
Z_i = H^1 \bigl(U,\, \sN_i^* \bigr) \subsetneq H^1 \bigl(U,\, \Omega^1_U \bigr).
\]
But since $c_1(\sL)$ is proportional to any of the $c_1(\sN_\sbullet)$, we find that
$c_1(\sL)$ is contained in $Z_1 \cap\cdots\cap Z_n = \{0\}$. The assertion is thus shown.
\end{proof}

Applying Proposition~\ref{prop:druel} to the smooth locus of a klt variety, we
obtain the following characterisation of quasi-abelian varieties.

\begin{thm}\label{thm:druel}
Let $X$ be a normal, projective klt variety of dimension $n \geq 2$. Assume that
the sheaf $\Omega^{[1]}_X$ of reflexive differentials is of the form
$\Omega^{[1]}_X \cong \sL^{\oplus n}$, where $\sL$ is reflexive of rank one. Then, $X$ is
quasi-abelian.
\end{thm}
\begin{proof}
We have seen in Proposition~\ref{prop:druel} that the $\bQ$-Cartier sheaves $\sL$
are numerically trivial, and then so is $\omega_X \cong \sL^{[\otimes n]}$. Since $X$ is klt,
it follows from \cite[Prop.\,V.4.9]{Nakayama04} or \cite[Th.\,4.2]{MR2134273}
that $\omega_X$ and $\sL$ are torsion. As a consequence, we find a finite
quasi-étale cover $\gamma: \what{X} \to X$ such that the reflexive pull-back
$\gamma^{[*]} \sL$ is trivial. The space $\what{X}$ is klt and its tangent sheaf
$\sT_{\what{X}} \cong \left(\gamma^{[*]} \sL \right)^{\oplus n}$ is likewise trivial. By the
solution of the Lipman-Zariski conjecture for spaces with klt singularities,
\cite[Th.\,6]{GKKP11}, it follows that the covering space $\what{X}$ is
smooth. In other words, $\what{X}$ is a parallelisable projective manifold.
By a classic result of Wang, $\what{X}$ is an Abelian variety; see for example
\cite[\S 3.9, Lem.\,1]{MR1334091}.
\end{proof}

\subsection{Proof of Proposition~\ref*{prop:icq}}
\label{ssec:popicq}

The proof of Proposition~\ref{prop:icq} uses the following lemma that might be
of independent interest.

\begin{lem}\label{lem:glcs}
Let $X$ be a normal, complex space with only cyclic quotient singularities.
Let $S \subseteq X_{\sing}$ be any irreducible component of the singular locus. Then,
there exists a non-empty, open set $S^\circ \subseteq S$ such that every $x \in S^\circ$ admits an
open neighbourhood $U=U(x) \subseteq X$ of the form $U = W \times V$, where $W$ is smooth
and $V$ has an isolated quotient singularity.
\end{lem}
\begin{proof}
To begin, consider the special case where $V$ is a complex vector space,
$\rho : \bZ_m \to \GL(V)$ is a linear representation of a cyclic group and where
$X \subseteq V / \bZ_m$ is an open neighbourhood of $[\vec 0]$. Recall that $\rho$
decomposes into direct a sum of one-dimensional representations. Identifying
$\bZ_m$ with roots of unity, we can thus find linear coordinates on $V$ such
that multiplication with $\xi \in \bZ_m$ is given as
\[
\xi\cdot (v_1, \dots, v_m) = \bigl(\xi^{n_1}\cdot v_1, \dots, \xi^{n_m}\cdot v_m\bigr).
\]
An elementary computation in these coordinates shows that $S^\circ$ exists, and can
in fact be chosen to be dense in $S$.

Returning to the general case where $X$ is arbitrary, let $s \in S$ be any
point. By~assumption, there exists an open neighbourhood $U = U(x)$ and a
cyclic cover $\gamma : \wtilde U \to U$, say with group $G$. Shrinking $U$ and
passing to a subgroup, if need be, we may assume without loss of generality
that $\gamma$ is totally branched over $s$, that is $\gamma^{-1}(s) = \{ \wtilde s \}$.
The cyclic group thus acts on $\wtilde U$ and fixes $\wtilde s$. The group
$G$ clearly acts on the vector space $T_{\wtilde U}|_{\wtilde s}$ and
linearisation at the fixed point, \cite[Proof of Th.\,4]{MR0084174} or
\cite[Cor.\,2 on p.\,13]{MR782881}, shows the existence of a $G$-invariant open
neighbourhood $\wtilde U'$ of $\vec 0 \in T_{\wtilde U}|_{\wtilde s}$ and an
equivariant, biholomorphic map between $\wtilde U$ and $\wtilde U$. In other
words, we are in the situation discussed in the first paragraph of this proof,
where the claim has already been shown.
\end{proof}

\begin{proof}[Proof of Proposition~\ref{prop:icq}]
As a first step in the proof of Proposition~\ref{prop:icq}, we claim that $\sL$
is $\bQ$-Cartier. In fact, taking determinants of both sides in \eqref{eq:icq},
we obtain that $\sL^{[\otimes \dim X]} \cong \omega_X$, which is $\bQ$-Cartier by assumption.
Choose a minimal number $N \in \bN^+$ such that $\sL^{[\otimes N]}$ is locally free. The
reflexive tensor product $\bigl(\Omega^{[1]}_U \bigr)^{[\otimes N]}$ is then likewise
locally free.

As a second step, we show that $X$ has cyclic quotient singularities. In
fact, given any $x \in X$, we find an open neighbourhood $U=U(x)$ over which
$\sL^{[\otimes N]}$ is trivial. Chose a trivialisation and let $\gamma: \wtilde{U} \to U$ be
the associated index-one cover, which is cyclic of order $N$. The complex
space $\wtilde{U}$ has again klt singularities and
$\gamma^{[*]}\, \sL \cong\nobreak \sO_{\wtilde{U}}$. In particular, it follows that
$\Omega^{[1]}_{\wtilde{U}} \cong \gamma^{[*]} \Omega^{[1]}_U$ and $\sT_{\wtilde U}$ are both free.
The solution of the Lipman-Zariski conjecture for spaces with klt
singularities, \cite[Th.\,3.8 \& comments after Th.\,1.1]{Druel13a} or
\cite[Th.\,1.13]{KS18}, then asserts that $\wtilde{U}$ is smooth. We conclude
that $X$ has cyclic quotient singularities only.

As a third and last step, we show that the singularities of $X$ are isolated.
Assume to the contrary and let $S \subseteq X_{\sing}$ be a positive-dimensional,
irreducible component. We have seen in Lemma~\ref{lem:glcs} that there a
non-empty, open subset $S^\circ \subseteq S$ such that every $x \in S^\circ$ admits an open
neighbourhood $U = U(x) \subseteq X$ of product form $U = W \times V$, where $W$ is smooth
and $V$ has an isolated quotient singularity. In particular,
$\dim W = \dim S$ and
\[
\Omega^{[1]}_U \cong \Omega^1_W \oplus \Omega^{[1]}_V.
\]
In particular, the reflexive tensor power $\bigl(\Omega^{[1]}_U \bigr)^{[\otimes N]}$ is
written as a direct sum of reflexive sheaves with
$\bigl(\Omega^{[1]}_W \bigr)^{[\otimes (N-1)]} \otimes \Omega^{[1]}_V$ as one of its direct
summands. It follows that $\bigl(\Omega^{[1]}_U \bigr)^{[\otimes N]}$ is not locally
free and accordingly, neither is $\bigl(\Omega^{[1]}_X \bigr)^{[\otimes N]}$. This
contradicts the results obtained in the first paragraph of this proof, and
therefore finishes the proof of Proposition~\ref{prop:icq}.
\end{proof}

% End literal retained input author-body.tex

\begin{thebibliography}{99}
\bibitem{MR1334091} D. N. Akhiezer - Lie group actions in complex analysis , Aspects of Mathematics , vol. E27 , Friedr. Vieweg \& Sohn, Braunschweig, 1995
\bibitem{MR925989} R. C. Alperin - “An elementary account of Selberg’s lemma” , Enseign. Math. (2) 33 (1987) no. 3-4, p. 269-273
\bibitem{MR2134273} F. Ambro - “The moduli $b$-divisor of an lc-trivial fibration” , Compositio Math. 141 (2005) no. 2, p. 385-403
\bibitem{MR670921} V. Ancona - “Faisceaux amples sur les espaces analytiques” , Trans. Amer. Math. Soc. 274 (1982) no. 1, p. 89-100
\bibitem{MR0261635} P. F. Baum \& R. Bott - “On the zeros of meromorphic vector-fields” , in Essays on Topology and Related Topics (Mémoires dédiés à Georges de Rham) , Springer, New York, 1970, p. 29-47
\bibitem{MR3286535} C. Birkar \& J. A. Chen - “Varieties fibred over abelian varieties with fibres of log general type” , Adv. Math. 270 (2015), p. 206-222
\bibitem{Bea83} A. Beauville - “Variétés kähleriennes dont la première classe de Chern est nulle” , J. Differential Geom. 18 (1983) no. 4, p. 755-782
\bibitem{MR1760872} A. Beauville - “Complex manifolds with split tangent bundle” , in Complex analysis and algebraic geometry , de Gruyter, Berlin, 2000, p. 61-70
\bibitem{HBPV} W. P. Barth, K. Hulek, C. A. M. Peters \& A. Van de Ven - Compact complex surfaces , Ergeb. Math. Grenzgeb. (3) , vol. 4 , Springer-Verlag, Berlin, 2004
\bibitem{BS76} C. Bănică \& O. Stănăşilă - Algebraic methods in the global theory of complex spaces , Editura Academiei, Bucharest; John Wiley \& Sons, London-New York-Sydney, 1976
\bibitem{Campana91} F. Campana - “On twistor spaces of the class $\mathscr{C}$” , J. Differential Geom. 33 (1991) no. 2, p. 541-549
\bibitem{MR0084174} H. Cartan - “Quotient d’un espace analytique par un groupe d’automorphismes” , in Algebraic geometry and topology. A symposium in honor of S. Lefschetz , Princeton University Press, Princeton, NJ, 1957, p. 90-102
\bibitem{MR3314830} F. Campana, B. Claudon \& P. Eyssidieux - “Représentations linéaires des groupes kählériens: factorisations et conjecture de Shafarevich linéaire” , Compositio Math. 151 (2015) no. 2, p. 351-376
\bibitem{CKP12} F. Campana, V. Koziarz \& M. Păun - “Numerical character of the effectivity of adjoint line bundles” , Ann. Inst. Fourier (Grenoble) 62 (2012) no. 1, p. 107-119
\bibitem{Debarre01} O. Debarre - Higher-dimensional algebraic geometry , Universitext , Springer-Verlag, New York, 2001
\bibitem{Dem02} J.-P. Demailly - “On the Frobenius integrability of certain holomorphic $p$-forms” , in Complex geometry (Göttingen, 2000) , Springer, Berlin, 2002, p. 93-98
\bibitem{MR1758581} S. Druel - “Variétés algébriques dont le fibré tangent est totalement décomposé” , J. reine angew. Math. 522 (2000), p. 161-171
\bibitem{Druel13a} S. Druel - “The Zariski-Lipman conjecture for log canonical spaces” , Bull. London Math. Soc. 46 (2014) no. 4, p. 827-835
\bibitem{MR2944479} O. Fujino \& Y. Gongyo - “On canonical bundle formulas and subadjunctions” , Michigan Math. J. 61 (2012) no. 2, p. 255-264
\bibitem{FL81} W. Fulton \& R. Lazarsfeld - “Connectivity and its applications in algebraic geometry” , in Algebraic geometry (Chicago, Ill., 1980) , Lect. Notes in Math. , vol. 862 , Springer, Berlin, 1981, p. 26-92
\bibitem{Fuj13} O. Fujino - “On maximal Albanese dimensional varieties” , Proc. Japan Acad. Ser. A Math. Sci. 89 (2013) no. 8, p. 92-95
\bibitem{GKK08} D. Greb, S. Kebekus \& S. J. Kovács - “Extension theorems for differential forms, and Bogomolov-Sommese vanishing on log canonical varieties” , Compositio Math. 146 (2010), p. 193-219, A slightly extended version is available as arXiv:0808.3647
\bibitem{GKKP11} D. Greb, S. Kebekus, S. J. Kovács \& T. Peternell - “Differential forms on log canonical spaces” , Publ. Math. Inst. Hautes Études Sci. 114 (2011) no. 1, p. 87-169, An extended version with additional graphics is available as arXiv:1003.2913
\bibitem{GKP13} D. Greb, S. Kebekus \& T. Peternell - “Étale fundamental groups of Kawamata log terminal spaces, flat sheaves, and quotients of abelian varieties” , Duke Math. J. 165 (2016) no. 10, p. 1965-2004
\bibitem{GKP20a} D. Greb, S. Kebekus \& T. Peternell - “Projective flatness over klt spaces and uniformisation of varieties with nef anti-canonical divisor” , 2020
\bibitem{GKPT19b} D. Greb, S. Kebekus, T. Peternell \& B. Taji - “The Miyaoka-Yau inequality and uniformisation of canonical models” , Ann. Sci. École Norm. Sup. (4) 52 (2019) no. 6, p. 1487-1535
\bibitem{GKPT17} D. Greb, S. Kebekus, T. Peternell \& B. Taji - “Nonabelian Hodge theory for klt spaces and descent theorems for vector bundles” , Compositio Math. 155 (2019) no. 2, p. 289-323
\bibitem{GKPT19} D. Greb, S. Kebekus, T. Peternell \& B. Taji - “Harmonic metrics on Higgs sheaves and uniformization of varieties of general type” , Math. Ann. 378 (2020) no. 3-4, p. 1061-1094
\bibitem{GKT16} D. Greb, S. Kebekus \& B. Taji - “Uniformisation of higher-dimensional varieties” , in Algebraic Geometry (Salt Lake City, 2015) (T. de Fernex, B. Hassett, M. Mustaţă, M. Olsson, M. Popa \& R. Thomas, eds.) , Proc. Sympos. Pure Math. , vol. 97 , American Mathematical Society, Providence, RI, 2018, p. 277-308
\bibitem{GrCa60} A. Grothendieck - “Techniques de construction en géométrie analytique. V. Fibrés vectoriels, fibrés projectifs, fibrés en drapeaux” , in Séminaire Henri Cartan , vol. 13, no. 1 , Secrétariat mathématique, Paris, 1960–1961, p. 1-15, Exp. no. 12
\bibitem{Ha77} R. Hartshorne - Algebraic geometry , Graduate Texts in Math. , vol. 52 , Springer-Verlag, New York, 1977
\bibitem{HMcK07} C. D. Hacon \& J. McKernan - “On Shokurov’s rational connectedness conjecture” , Duke Math. J. 138 (2007) no. 1, p. 119-136
\bibitem{MR782881} A. Huckleberry \& E. Oeljeklaus - Classification theorems for almost homogeneous spaces , vol. 9 , Université de Nancy, Institut Élie Cartan, Nancy, 1984
\bibitem{MR3030068} P. Jahnke \& I. Radloff - “Semistability of restricted tangent bundles and a question of I. Biswas” , Internat. J. Math. 24 (2013) no. 1, article ID 1250122, 15 pages
\bibitem{Kawamata85} Y. Kawamata - “Minimal models and the Kodaira dimension of algebraic fiber spaces” , J. reine angew. Math. 363 (1985), p. 1-46
\bibitem{MR782236} Y. Kawamata - “Pluricanonical systems on minimal algebraic varieties” , Invent. Math. 79 (1985) no. 3, p. 567-588
\bibitem{KM98} J. Kollár \& S. Mori - Birational geometry of algebraic varieties , Cambridge Tracts in Mathematics , vol. 134 , Cambridge University Press, Cambridge, 1998
\bibitem{KMM87} Y. Kawamata, K. Matsuda \& K. Matsuki - “Introduction to the minimal model problem” , in Algebraic geometry (Sendai, 1985) , Adv. Stud. Pure Math. , vol. 10 , North-Holland, Amsterdam, 1987, p. 283-360
\bibitem{Kob87} S. Kobayashi - Differential geometry of complex vector bundles , Publications of the Math. Society of Japan , vol. 15 , Iwanami Shoten and Princeton University Press, Princeton, NJ, 1987
\bibitem{Kollar93} J. Kollár - “Shafarevich maps and plurigenera of algebraic varieties” , Invent. Math. 113 (1993) no. 1, p. 177-215
\bibitem{MR1341589} J. Kollár - Shafarevich maps and automorphic forms , M. B. Porter Lectures , Princeton University Press, Princeton, NJ, 1995
\bibitem{KS18} S. Kebekus \& C. Schnell - “Extending holomorphic forms from the regular locus of a complex space to a resolution of singularities” , J. Amer. Math. Soc. 34 (2021) no. 2, p. 315-–368
\bibitem{Lai11} C.-J. Lai - “Varieties fibered by good minimal models” , Math. Ann. 350 (2011) no. 3, p. 533-547
\bibitem{LT18} S. Lu \& B. Taji - “A characterization of finite quotients of abelian varieties” , Internat. Math. Res. Notices (2018) no. 1, p. 292-319
\bibitem{NakayamaNormalizedPreprint} N. Nakayama - “Normalized tautological divisors of semi-stable vector bundles” (1998), RIMS preprint 1214, available at https://www.kurims.kyoto-u.ac.jp/preprint/preprint\_y1998.html
\bibitem{Nakayama04} N. Nakayama - Zariski-decomposition and abundance , MSJ Memoirs , vol. 14 , Mathematical Society of Japan, Tokyo, 2004
\bibitem{MR0260758} M. Raynaud - Faisceaux amples sur les schémas en groupes et les espaces homogènes , Lect. Notes in Math. , vol. 119 , Springer-Verlag, Berlin-New York, 1970
\bibitem{Rossi68} H. Rossi - “Picard variety of an isolated singular point” , Rice Univ. Studies 54 (1968) no. 4, p. 63-73
\bibitem{Schwald16} M. Schwald - “Low degree Hodge theory for klt varieties” , 2016
\bibitem{MR0287033} Y.-T. Siu \& G. Trautmann - Gap-sheaves and extension of coherent analytic subsheaves , Lect. Notes in Math. , vol. 172 , Springer-Verlag, Berlin, 1971
\bibitem{Takayama2003} S. Takayama - “Local simple connectedness of resolutions of log-terminal singularities” , Internat. J. Math. 14 (2003) no. 8, p. 825-836
\bibitem{MR641815} E. Viehweg - “Die Additivität der Kodaira Dimension für projektive Faserräume über Varietäten des allgemeinen Typs” , J. reine angew. Math. 330 (1982), p. 132-142
\end{thebibliography}
\end{document}
